\documentclass[10pt,a4paper]{amsart}
\usepackage[margin=19mm]{geometry}
\usepackage{amsmath,amssymb,mathtools}
\usepackage[scr=rsfs,cal=euler]{mathalfa}
\usepackage{xfrac}
\usepackage[unicode,hidelinks,bookmarks=false]{hyperref}
\newcommand{\ie}{i.e.\ }
\newcommand{\cf}{cf.\ }
\newcommand{\resp}{respectively\ }
\newcommand{\Subsection}{\subsection}
\providecommand{\nobreakdash}{\nobreak\hbox{-}}
\setlength{\emergencystretch}{2em}
\multlinegap0pt
\usepackage{bm}
\usepackage{bbm}
\usepackage{dsfont}
\usepackage{xspace}
\usepackage{cancel}
\usepackage{mathtools}
\usepackage{amsthm}
\makeatletter
\@ifundefined{theorem}{\newtheorem{theorem}{Theorem}}{}
\@ifundefined{lemma}{\newtheorem{lemma}[theorem]{Lemma}}{}
\@ifundefined{proposition}{\newtheorem{proposition}[theorem]{Proposition}}{}
\makeatother
\makeatletter
\newcommand{\Ls}{\mathcal{L}}
\newcommand{\dif}{{\mathrm{d}}}
\def\eqref#1{(\ref{#1})}
\def\gT{{\mathcal{T}}}
\def\mH{{\bm{H}}}
\def\mJ{{\bm{J}}}
\def\mK{{\bm{K}}}
\def\mP{{\bm{P}}}
\expandafter\ifx\csname manualassumption\endcsname\relax
\newenvironment{manualassumption}[1]{%
  \renewcommand\themanualtheoreminner{#1}%
  \manualtheoreminner
}{\endmanualtheoreminner}
\fi
\def\vw{{\bm{w}}}
\def\vx{{\bm{x}}}
\def\vy{{\bm{y}}}
\makeatother
\title[Solving Inverse Problems via Diffusion-Based Priors: An Appr]{Solving Inverse Problems via Diffusion-Based Priors: An Approximation-Free Ensemble Sampling Approach}
\author{Haoxuan Chen, Yinuo Ren, Martin Renqiang Min, Lexing Ying, Zachary Izzo}
\date{Source: arXiv:2506.03979v2 (CC BY 4.0)}
\begin{document}
\maketitle

\noindent\emph{Source and licence.} Solving Inverse Problems via Diffusion-Based Priors: An Approximation-Free Ensemble Sampling Approach. Version arXiv:2506.03979v2 (CC BY 4.0), distributed under the Creative Commons Attribution 4.0 International licence, as recorded in the arXiv licence metadata for this exact version and shown on the abstract page https://arxiv.org/abs/2506.03979v2. Full rights notice: licenses/arxiv-2506.03979-CC-BY-4.0.txt.\par\smallskip

\noindent\textit{Editorial scope.} Counted unit: the Lemma whose source label is \texttt{lem: derivation of single particle dynamics} in the article (source0.tex lines 1562--1590, source label \texttt{lem: derivation of single particle dynamics}), delivered together with the article's own complete original proof of it (source0.tex lines 1710--1779, one \texttt{proof} environment of 70 lines). No mathematics is written, repaired or re-derived: statement and proof are the article's own text line for line. The proof is detached from the statement in the source: the article states the result at lines 1562--1590 and prints its proof at lines 1710--1779, and the proof environment's own header names the statement, so the association is the article's own. Required context retained uncounted: eqn: weak form of joint PDE (lines 1641--1648); eqn: intermediate derivaion of weak form for joint PDE (lines 1663--1671); eqn: notation for the sde drifts (lines 1679--1686); eqn: derivation of weak form for joint density (lines 1695--1703). Every definition, display and label of those lines is retained unchanged. Omitted from this package and not redistributed: 1--1561, 1591--1640, 1649--1662, 1672--1678, 1687--1694, 1704--1709, 1780--2612. Nothing inside a retained excerpt is omitted or altered: every retained body is delivered exactly as the article's lines are written, so no omission receipt of any kind accompanies this package. Citations: the retained excerpts carry no citation command at all, so no bibliography entry of the article is transcribed and no reference is redistributed. Reference adaptation: none was needed and none was made; every reference inside the delivered unit resolves inside it, so no unresolved reference or citation is produced. Printed slips kept literal: no printed word, symbol, hypothesis or number of the counted statement or of its proof was altered. Layout and class: the article's own class is \texttt{article} with packages including natbib, neurips\_2025, inputenc, fontenc, hyperref, url, booktabs, amsfonts, nicefrac, microtype, xcolor, graphicx, wrapfig, subcaption; the wrapper is the lane's pinned \texttt{amsart} editorial wrapper at 10pt on A4, which restates the theorem-like environments the retained text needs and adds the article's own macro definitions that the retained text needs (\texttt{\textbackslash Ls}, \texttt{\textbackslash dif}, \texttt{\textbackslash eqref}, \texttt{\textbackslash gT}, \texttt{\textbackslash mH}, \texttt{\textbackslash mJ}, \texttt{\textbackslash mK}, \texttt{\textbackslash mP}, \texttt{\textbackslash vw}, \texttt{\textbackslash vx}, \texttt{\textbackslash vy}), with extra wrapper packages (bm, bbm, dsfont, xspace, cancel, mathtools). The delivered numbering is the unit's own. Assets: no retained span contains a figure environment, a graphics-inclusion command, a PDF-page inclusion, a table or a TikZ picture; no image file is copied into this package, the package declares no asset dependency and no image file is redistributed with it. Licence: version arXiv:2506.03979v2 of this article is distributed under the Creative Commons Attribution 4.0 International licence, as recorded in the arXiv licence metadata for this exact version and shown on the abstract page https://arxiv.org/abs/2506.03979v2. A full rights notice is delivered at licenses/arxiv-2506.03979-CC-BY-4.0.txt. Characters: every retained excerpt is pure ASCII, so the delivered engine, XeLaTeX with its default Latin Modern text font, reports no missing character.
\begin{proposition}
    For any test function $\varphi: \mathbb{R}^n \times \mathbb{R} \rightarrow \mathbb{R}$, the joint distribution $\gamma_t = \gamma_t(\vx, \beta)$ satisfies the following PDE:
    \begin{equation}
        \frac{\partial}{\partial t}\gamma_t = -\nabla_{\vx} \cdot \left(\mK_{t}\gamma_t\right) - \frac{\partial}{\partial \beta}\left(b_t \gamma_t\right)+\frac{1}{2}V(t)^2\Delta_{\vx}\gamma_t,
        \label{eqn: weak form of joint PDE}
    \end{equation}
    with the initial condition $\gamma_0(\vx,\beta) =\widehat{q}_{\vy}(\vx,0) \times \delta_{\beta = 1}$.
\end{proposition}

    \begin{equation}
    \label{eqn: intermediate derivaion of weak form for joint PDE}
    \begin{aligned}
    &\left\langle\varphi, \frac{\partial}{\partial t}\gamma_t\right\rangle = \frac{\dif}{\dif t}\left\langle\varphi,\gamma_t\right\rangle = \frac{\dif}{\dif t}\left\langle\gT^{(\vx,\beta)}_{t}\varphi, \gamma_0\right\rangle \\
    =& \left\langle\frac{\partial}{\partial t}gT^{(\vx,\beta)}_{t}\varphi, \gamma_0\right\rangle 
    = \left\langle\gT^{(\vx,\beta)}_{t}\circ\Ls^{(\vx,\beta)}\varphi, \gamma_0\right\rangle 
    = \left\langle\Ls^{(\vx,\beta)}\varphi, \gamma_t\right\rangle.  
    \end{aligned}    
    \end{equation}

    \begin{equation}
    \label{eqn: notation for the sde drifts}
    \begin{aligned}
    \mK_t(\vx) &= \widehat{\mH}(\vx,t)-V(t)^2\nabla_{\vx}\mu_{\vy}(\vx),\\
    b_t(\vx,\beta) &= \left(U(\vx,t) - \widehat{\mH}(\vx,t)^\intercal\nabla_{\vx}\mu_{\vy}(\vx)\right)\beta\\
    &- \left(\int_{\mathbb{R}^n}\left(U(\vx^\ast,t) - \widehat{\mH}(\vx^\ast,t)^\intercal\nabla_{\vx}\mu_{\vy}(\vx^\ast)\right)\left(P_{\beta}\gamma_t\right)(\vx^\ast)\dif \vx^\ast\right)\beta.
    \end{aligned}
    \end{equation}

    \begin{equation}
    \label{eqn: derivation of weak form for joint density}
    \begin{aligned}
    &\left\langle\Ls^{(\vx,\beta)}\varphi, \gamma_t\right\rangle = \left\langle\left(\nabla_{\vx}\varphi\right)^\intercal\mK_t + \frac{\partial \varphi}{\partial \beta}b_t + \frac{1}{2}V(t)^2\Delta_{\vx}\varphi, \gamma_t\right\rangle\\
    =& \sum_{i=1}^{n}\left\langle\frac{\partial \varphi}{\partial x_i},\mK_{t,i} \gamma_{t} \right\rangle + \left\langle\frac{\partial \varphi}{\partial \beta},b_t\gamma_{t} \right\rangle + \frac{1}{2}V(t)^2\sum_{i=1}^{n}\left\langle\frac{\partial^2 \varphi}{\partial x_i^2}, \gamma_{t} \right\rangle\\
    =& -\left\langle\varphi, \sum_{i=1}^{n}\frac{\partial}{\partial x_i}\left(\mK_{t,i} \gamma_{t}\right) + \frac{\partial}{\partial \beta}\left(b_t \gamma_{t}\right) + \frac{1}{2}V(t)^2\Delta_{\vx}\gamma_{t} \right\rangle\\
    =& \left\langle\varphi, -\nabla_{\vx} \cdot \left(\mK_{t}\gamma_t\right) - \frac{\partial}{\partial \beta}\left(b_t \gamma_t\right)+\frac{1}{2}V(t)^2\Delta_{\vx}\gamma_t\right\rangle,
    \end{aligned}    
    \end{equation}

\begin{lemma}
    \label{lem: derivation of single particle dynamics}    
    Consider a single particle $(\vx_t,\beta_t)$ governed by
    \begin{equation}
    \label{eqn: PDE soln via single particle app} 
    \begin{cases}
    \dif \vx_t &= \left(\widehat{\mH}(\vx_t,t)-V(t)^2\nabla_{\vx}\mu_{\vy}(\vx_t)\right)\dif t + V(t)\dif \vw_t,\\
    \dif \beta_t &= \left(U(\vx_t,t) - \widehat{\mH}(\vx_t,t)^\intercal\nabla_{\vx}\mu_{\vy}(\vx_t)\right)\beta_t \dif t \\
    &- \left(\displaystyle\int_{\mathbb{R}^n}\left(U(\vx,t) - \widehat{\mH}(\vx,t)^\intercal\nabla_{\vx}\mu_{\vy}(\vx)\right)\left(P_{\beta}\gamma_t\right)(\vx)\dif \vx\right)\beta_t \dif t,    
    \end{cases}
    \end{equation}
    with initial condition $\vx_0 = \vx^\ast$ and $\beta_0 = 1$, where $\vx^\ast$ is sampled from the initial posterior distribution $\widehat{q}_{\vy}(\vx,0)$, $(\vw_t)_{t\geq 0}$ is a standard Brownian motion in $\mathbb{R}^n$, $\gamma_t(\vx,\beta)$ denotes the joint probability distribution of $(\vx_t,\beta_t)$ on $\mathbb{R}^n \times \mathbb{R}$, 
    $$
        P_{\beta}\gamma_t(\vx) := \int_{\mathbb{R}}\beta\gamma_{t}(\vx,\beta)\dif\beta
    $$ 
    denotes the weighted projection of $\gamma_t$ onto $\vx$, and 
    $$
        U(\vx,t) := \frac{1}{2}V(t)^2\left(\|\nabla_{\vx}\mu_{\vy}(\vx)\|_2^2-\Delta_{\vx}\mu_{\vy}(\vx)\right).
    $$ 
    Then we have that $P_{\beta}\gamma_t(\vx) = \widehat{q}_{\vy}(\vx,t)$ for any $\vx \in \mathbb{R}^n$ and $t \in [0,T]$, \emph{i.e.} $P_{\beta}\gamma_t(\cdot)$ solves the following PDE:
    \begin{equation}
            \label{eqn: PDE of posterior diffusion app2}
            \begin{aligned}
            \frac{\partial}{\partial t}&\widehat{q}_{\vy} = -\nabla_{\vx} \cdot \left(\left(\widehat{\mH}(\vx,t)-V(t)^2\nabla_{\vx}\mu_{\vy}\right)\widehat{q}_{\vy}\right)
            +\frac{1}{2}V(t)^2\Delta_{x}\widehat{q}_{\vy}\\
            &+\left(U(\vx,t) - \widehat{\mH}(\vx,t)^\intercal\nabla_{\vx}\mu_{\vy} - \int_{\mathbb{R}^n}\left(U(\vx,t) - \widehat{\mH}(\vx,t)^\intercal\nabla_{\vx}\mu_{\vy}\right)\widehat{q}_{\vy}\dif \vx\right)\widehat{q}_{\vy}.
            \end{aligned}
        \end{equation}
\end{lemma}

\begin{proof}[Proof of Lemma~\ref{lem: derivation of single particle dynamics}]
    
    
    By defining 
    $$\gamma^{\mP}_t(\vx) := P_{\beta}\gamma_t(\vx) = \int_{\mathbb{R}}\beta\gamma_{t}(\vx,\beta)\dif\beta$$
    to be the weighted projection of $\gamma_t$, we then have that $\gamma^{\mP}_0(\vx) = \widehat{q}_{\vy}(\vx,0)$. 
    Below we proceed to derive the PDE govering the evolution of $\gamma^{\mP}_t$ based on~\eqref{eqn: weak form of joint PDE}. 
    
    For any test function $\psi:\mathbb{R}^n \rightarrow \mathbb{R}$, taking $\varphi(\vx,\beta) = \beta\psi(\vx):\mathbb{R}^{n} \times \mathbb{R} \rightarrow \mathbb{R}$ in the weak form derived in~\eqref{eqn: intermediate derivaion of weak form for joint PDE} and~\eqref{eqn: derivation of weak form for joint density} yields
    \begin{equation}
    \label{eqn: intermediate weak form for PDE of projected measure}
    \begin{aligned}
    &\left\langle\psi, \frac{\partial}{\partial t}\gamma^{\mP}_t\right\rangle = \frac{\dif}{\dif t}\left\langle\psi,\gamma^{\mP}_t\right\rangle= \frac{\dif}{\dif t}\left(\int_{\mathbb{R}^n}\psi(\vx)\left(\int_{\mathbb{R}}\beta\gamma_{t}(\vx,\beta)\dif\beta\right)\dif\vx\right) \\
    =& \frac{\dif}{\dif t}\left\langle\varphi,\gamma_t\right\rangle = \left\langle\varphi, \frac{\partial}{\partial t}\gamma_t\right\rangle
    = \left\langle\varphi, -\nabla_{\vx} \cdot \left(\mK_{t}\gamma_t\right) - \frac{\partial}{\partial \beta}\left(b_t \gamma_t\right)+\frac{1}{2}V(t)^2\Delta_{\vx}\gamma_t\right\rangle\\
    =& -\left\langle \varphi, \nabla_{\vx} \cdot \left(\mK_{t}\gamma_t\right) \right\rangle -\left\langle \varphi, \frac{\partial}{\partial \beta}\left(b_t \gamma_t\right) \right\rangle + \frac{1}{2}V(t)^2\left\langle \varphi, \Delta_{\vx}\gamma_t \right\rangle.
    \end{aligned}    
    \end{equation}

    For the first and third terms in the last expression of~\eqref{eqn: intermediate weak form for PDE of projected measure}, we can further simplify them as follows
    \begin{equation}
    \label{eqn: derivation of first term}
    \begin{aligned}
    \left\langle \varphi, \nabla_{\vx} \cdot \left(\mK_{t}\gamma_t\right) \right\rangle &= \int_{\mathbb{R}}\beta\left(\int_{\mathbb{R}^n}\psi(\vx)\left(\nabla_{\vx} \cdot\left(\mK_t(\vx) \gamma_t(\vx,\beta)\right)\right)\dif\vx\right)\dif\beta\\
    &= \int_{\mathbb{R}^n}\psi(\vx)\left(\nabla_{\vx} \cdot\left(\mK_t(\vx) \left(\int_{\mathbb{R}}\beta\gamma_t(\vx,\beta)\dif\beta\right)\right)\right)\dif\vx\\
    &= \int_{\mathbb{R}^n}\psi(\vx)\left(\nabla_{\vx} \cdot\left(\mK_t(\vx)\gamma^{\mP}_t(\vx)\right)\right)\dif\vx \\
    &=\left\langle\psi,\nabla_{\vx}\cdot(\mK_{t}\gamma^{\mP}_t)\right\rangle
    \end{aligned}    
    \end{equation}
    and 
    \begin{equation}
        \label{eqn: derivation of third term}
        \begin{aligned}
            \left\langle \varphi, \Delta_{\vx}\gamma_t \right\rangle &= \int_{\mathbb{R}}\beta\left(\int_{\mathbb{R}^n}\psi(\vx)\left(\Delta_{\vx}\gamma_t(\vx,\beta)\right)\dif\vx\right)\dif\beta\\
            &= \int_{\mathbb{R}^n}\psi(\vx)\left(\Delta_{\vx}\left(\int_{\mathbb{R}}\beta\gamma_t(\vx,\beta)\dif\beta\right)\right)\dif\vx\\
            &= \int_{\mathbb{R}^n}\psi(\vx)\left(\Delta_{\vx}\gamma^{\mP}_t\right)(\vx)\dif\vx = \left\langle\psi,\Delta_{\vx}\gamma^{\mP}_t\right\rangle,
        \end{aligned}
    \end{equation}
    respectively.

    Moreover, for the second term in the last expression of~\eqref{eqn: intermediate weak form for PDE of projected measure}, we use 
    \begin{equation}
    \label{eqn: defn of J_t}
    \mJ_t(\vx) := U(\vx,t) - \widehat{\mH}(\vx,t)^\intercal\nabla_{\vx}\mu_{\vy}(\vx)
    \end{equation} 
    to denote the integrand in the definition of $b_t$, which helps us rewrite $b_t$ as follows
    $$b_t(\vx,\beta) = \left(\mJ_t(\vx)-\int_{\mathbb{R}^n}\mJ_t(\vx^\ast)\gamma^{\mP}_t(\vx^\ast)\dif\vx^\ast\right)\beta.$$

    Then we apply integration by parts again to deduce that
    \begin{equation}
    \label{eqn: derivation of reweighting term}
    \begin{aligned}
    &\left\langle \varphi, \frac{\partial}{\partial \beta}\left(b_t \gamma_t\right) \right\rangle = -\left\langle \frac{\partial}{\partial \beta}\varphi, b_t \gamma_t \right\rangle = -\left\langle \psi(\vx), b_t \gamma_t \right\rangle\\
    = &-\int_{\mathbb{R}}\psi(\vx)\gamma_t(\vx,\beta)\left(\mJ_t(\vx)-\int_{\mathbb{R}^n}\mJ_t(\vx^\ast)\gamma^{\mP}_t(\vx^\ast)\dif\vx^\ast\right)\beta\dif\vx\dif\beta\\
    = &-\int_{\mathbb{R}}\psi(\vx)\gamma^{\mP}_t(\vx)\left(\mJ_t(\vx)-\int_{\mathbb{R}^n}\mJ_t(\vx^\ast)\gamma^{\mP}_t(\vx^\ast)\dif\vx^\ast\right)\dif\vx\\
    = &-\left\langle\psi(\vx), \gamma^{\mP}_t(\vx)\left(\mJ_t(\vx)-\int_{\mathbb{R}^n}\mJ_t(\vx^\ast)\gamma^{\mP}_t(\vx^\ast)\dif\vx^\ast\right)\right\rangle.
    \end{aligned}    
    \end{equation}

    Substituting~\eqref{eqn: derivation of first term},~\eqref{eqn: derivation of third term}, and~\eqref{eqn: derivation of reweighting term} into~\eqref{eqn: intermediate weak form for PDE of projected measure} then gives us the weak form of the PDE governing the evolution of the projected measure $\gamma^{\mP}_t = P_{\beta}\gamma_t$. 
    
    Hence, we finally have that $\gamma^{\mP}_t(\vx) = P_{\beta}\gamma_t(\vx)$ satisfies the following PDE
    \begin{equation}
    \label{eqn: PDE of projected measure}
    \frac{\partial}{\partial t}\gamma^{\mP}_t = -\nabla_{\vx} \cdot \left(\mK_{t}\gamma^{\mP}_t\right) +\frac{1}{2}V(t)^2\Delta_{\vx}\gamma^{\mP}_t + \left(\mJ_t-\int_{\mathbb{R}^n}\mJ_t(\vx^\ast)\gamma^{\mP}_t(\vx^\ast)\dif\vx^\ast\right)\gamma^{\mP}_t,
    \end{equation}
    with initial condition $\gamma^{\mP}_0(\vx) =\widehat{q}_{\vy}(\vx,0)$. 
    
    Plugging in the expressions of $\mK_t$ and $\mJ_t$ given in~\eqref{eqn: notation for the sde drifts} and~\eqref{eqn: defn of J_t} indicates that equation~\eqref{eqn: PDE of projected measure} is exactly the PDE provided in~\eqref{eqn: PDE of posterior diffusion app2}. This concludes our proof.
\end{proof}

\end{document}
